\chapter{总结与展望}

\section{主要研究结论}

本文围绕胸部X光肺炎二分类完成了从数据审计到外部评价的完整课程设计。项目首先验证旧结果能否复算，随后发现Kermany官方train/test有170个肺炎subject ID交叉。这个问题使旧0.89--0.91准确率只能作为公开目录复现。项目重新建立患者级划分，冻结manifest并通过患者、哈希和标签冲突门禁。

针对RQ1，结论是患者级审计不可省略，但患者交叉不能被简单解释为旧分数虚高。新患者级随机划分内部更容易，三模型accuracy达到约0.98。新旧split同时改变了独立性和分布难度，必须分别讨论。

针对RQ2，DenseNet121、ConvNeXt-Tiny与ViT-B/16在患者独立内部test上的差异很小。三seed集成balanced accuracy为0.9727--0.9817，95\%区间重叠。没有证据支持某种结构在当前任务中具有明确优势。ViT配置漂移曾造成明显欠拟合，也说明合理训练配方比只写架构名称更重要。

针对RQ3，三模型直接迁移RSNA后accuracy降到约0.66--0.68，主要错误是阴性误报。简单图像统计的来源分类AUC为0.9551，证明两个来源存在明显风格差异。外部下降远大于模型结构差异，单一儿童胸片数据上的高分不能代表跨来源稳定性。

针对RQ4，低成本强度扰动和标签平滑把RSNA balanced accuracy提高到约0.7116，内部保持稳定，但外部AUC下降，只能称为部分改善。简单混合训练取得最完整结果：RSNA balanced accuracy约0.7645、AUC约0.8565，Kermany内部balanced accuracy约0.9684；域均衡波动更大且没有超过简单混合。混合收益依赖RSNA有监督数据，不等于对未知医院泛化。

校准分析进一步显示，Kermany-only模型外部温度平均约6.81，存在严重过度自信；简单混合外部温度约1.56，概率漂移较小。validation冻结阈值能改善RSNA操作点，但阈值跨seed不稳定。Grad-CAM和错误网格只能帮助解释模型行为，不能作为病灶定位。

工程方面，项目补齐完整性审计、患者级数据生成、多seed聚合、温度范围检查和冻结阈值工具；网页支持单图推理与概率展示。所有输出仅用于课程设计和研究训练。当前证据不支持临床诊断、真实病灶定位或患者决策。

最终研究判断是：本项目的主要提升来自实验可信度和数据覆盖，而非网络数量。患者级独立、外部标签语义、统计不确定性和概率可靠性共同决定模型结论。把这些环节讲清楚，比在同一测试集继续寻找更高单次分数更符合课程设计训练目标。


\section{患者级审计改变了什么}

最直接的改变是证据等级。旧Kermany指标在数学上可复算，却不满足患者独立性；新版把它降为官方目录复现，并建立新manifest。这个过程说明“bug”不只指代码异常，也包括数据单位与研究问题不一致。若目标是对未见患者分类，患者必须是划分单位。

审计没有支持一个简单故事：去除患者交叉后分数并未下降，反而升到约0.98。原因是新划分同时改变了test分布。官方test可能包含更强来源/子类型偏移，新随机患者划分更接近同来源训练数据。因此患者泄漏风险与数据集难度必须分开；本文不声称170个交叉患者导致旧分数虚高多少。

\section{模型结构与数据覆盖}

三模型内部差异不足1个百分点，外部RSNA下降约30个百分点。来源分类器在排除纵横比后仍有0.9551 AUC。两项证据共同说明，当前瓶颈主要是数据域，而不是CNN或Transformer结构。

简单混合训练把RSNA AUC提高到约0.8565、balanced accuracy提高到0.7645，同时内部balanced accuracy保持0.9684。它的收益来自目标域有监督数据覆盖，不能推广为对任意新医院的域泛化。若再到第三个未见来源，仍可能下降。真正的域泛化应在开发时不使用最终目标域标签，并在多个外部机构验证。

\section{低成本稳健方案的双面结果}

标签平滑与强度扰动减少外部过度自信，Brier和specificity改善，但AUC下降。若只看accuracy或balanced accuracy，会把它写成成功；若只看AUC，又会完全否定它。更准确的理解是：它改变了概率尺度和0.5操作点，使误报减少，但排序能力没有增强。

这一结果说明评价指标必须对应使用目的。筛查任务可能重视高recall下specificity；概率风险沟通重视校准；排序或后续阈值选择重视AUC。课程设计没有临床成本函数，因此不选唯一“最佳”指标，而报告完整权衡。

\section{域均衡为何未胜出}

域均衡让RSNA小来源在采样中被放大，理论上可以防止Kermany主导。但有放回抽样重复少量RSNA图像，且两个域标签语义并不完全一致。结果表现为更大seed波动、内部specificity不稳、外部不如简单混合。

不能据此断言域均衡普遍无效。当前实现只按来源平衡，不按来源-类别联合组；训练epoch固定，超参数也未针对该策略调节。结论限于本项目设置：在同预算和当前数据量下，没有证据支持它优于简单混合。

\section{内部有效性威胁}

第一，Kermany subject ID来自文件名规则，不是官方患者表。规则与公开命名高度一致，但仍可能把不同患者错误合并或把同一患者拆开。第二，Flash Attention提示非确定位算法，三seed缓解但没有达到位级复现。第三，模型学习率和epoch按架构设置，不是完全相同超参数；比较回答“合理配方下模型表现”，不是纯架构因果效应。第四，新locked test在本轮已被多个方案评价，后续继续迭代会使其逐渐失去锁定性。

\section{构念有效性威胁}

三个数据集的“肺炎/阴性”并非相同构念。Kermany目录标签来自儿童病例分类；RSNA target围绕可能肺炎框；NIH为报告挖掘。RSNA阴性可能包含其他异常，NIH Pneumonia可能与其他标签共存。把它们统一成二分类便于代码，却可能把标签差异误当作视觉域偏移。

resize到正方形会拉伸不同纵横比图像。ImageNet三通道归一化用于复制预训练输入，却不一定是胸片最优。水平翻转可能改变侧别。未来应比较保持纵横比padding、DICOM窗宽窗位和肺野裁剪，并用validation选择。

\section{外部有效性威胁}

RSNA只使用1707张镜像可用子集，不代表完整挑战集；NIH仅130张。项目没有独立医院前瞻性数据，没有按年龄、性别、设备或体位做亚组评价。Kermany儿童人群与RSNA成人占比较高的人群差异很大。当前结论只能推广到这些公开数据的特定处理版本。

混合方案使用了RSNA train/val，因此RSNA test是“已知目标域的留出评价”，不是完全外部验证。课程报告必须避免把它写成无监督域泛化。要证明对未知来源有效，需要至少一个未参与任何方法设计的足量第三数据集。

\section{统计结论威胁}

三个seed只能初步估计训练波动。患者bootstrap处理了同一患者多图相关，但subject解析误差会影响区间。多个模型、指标和方案会产生多重比较；本文没有用大量显著性检验挑选结果，而依赖预设3个百分点门槛、双域约束和效果方向。

ECE依赖10个固定分箱，小样本和类别比例变化会影响数值。温度网格虽扩展到20并检查边界，仍是离散近似。validation阈值跨seed变化明显，提示单一阈值不稳定。

\section{主要工作与实践价值}

本项目的实际工作不只包括训练三个模型。更重要的部分是：发现并机器化检查患者交叉；修正配置漂移；建立患者级manifest；完成27个严格基线评价集合、18个稳健/混合核心训练-测试组合及校准、阈值和错误分析；把旧结果按证据等级重新分类；将网页与权重边界写清楚。

课程设计的价值体现在能解释为什么一个看似高分的模型仍可能不可靠，并用实验区分模型结构、数据覆盖、操作点和概率校准。负结果和停止规则也是工作成果，因为它们防止继续消耗算力追逐同一test最高值。

\section{研究展望}

第一，获取完整RSNA或另一个足量外部数据，预先冻结最终评价；第二，按patient、view和共病建立更严格标签；第三，比较保持纵横比、肺野分割和灰度标准化；第四，在新的外部集合上评价CORAL、GroupDRO等方法；第五，增加五seed和配对bootstrap差值区间；第六，由医学人员盲法复核固定抽样错误案例。

只有这些条件逐步满足后，才适合讨论临床研究设计。当前网页不应扩展成患者使用工具。

\section{关键问题讨论}

\subsection{为什么患者级新划分仍然高达98\%}

新划分消除了同一subject跨集合，却仍是同一公开数据来源的随机分层。训练和test在年龄、人群、设备与处理流程上接近，任务本身也可能有明显类别特征。高内部结果只说明同来源未见患者分类较容易，RSNA下降证明它不能推广成通用能力。

\subsection{既然官方test有患者交叉，旧结果是否全部作废}

不作废。预测和指标计算正确，也能复现公开目录benchmark；问题是不能解释为患者独立估计。报告保留旧表并明确降级。科学修正不等于删除不利历史，而是改变其能支撑的论断。

\subsection{为什么不把简单混合称为域泛化}

因为训练直接使用RSNA标签。它是有监督多域训练或目标域适配。域泛化通常要求最终目标域在开发时不可见。准确命名比方法听起来是否高级更重要。

\subsection{为什么不继续做更多模型和复杂算法}

三模型内部/外部区间重叠，结构不是主要矛盾。简单混合已满足预设门槛，继续在同一test挑方法会增加适应风险。剩余时间用于多seed、校准、阈值、错误审计和报告一致性，更能提高课程质量。

\subsection{论文内容如何避免依靠图像和附录扩充篇幅}

新版正文没有使用一页一图，主要内容由8章文字、公式和表格组成。正文详细记录数据单位、实验职责、多seed、置信区间、配置错误、负结果和有效性威胁。附录命令和产物索引位于参考文献之后，不计入正文目标。

\section{研究规范与工作边界}

项目使用自动化工具辅助代码审查、实验执行、数字汇总和文稿组织，但研究判断必须由可检查产物支撑。报告不把未执行的GroupDRO/CORAL写成完成，不把旧单seed结果包装成新版多seed，也不伪造医学人员复核。

所有自动生成文字都需要与JSON、CSV和日志核对。公开提交应保留代码、配置、manifest、审计报告和参考文献。若后续人工修改数字，必须重新运行一致性检查。团队成员应能说明各自负责的数据、模型、系统或写作内容，而不是只会复述最终摘要。

课程答辩中，最有说服力的不是宣称“导师挑不出问题”，而是主动指出患者交叉、标签差异和test适应，并展示如何修正。无法消除的限制应明确保留，这比过度承诺更符合研究训练目的。
